\subsubsection{Representing terms w/ variables in Twelf
using Higher-Order Abstract Syntax}
\label{sec:twelf-hoas}

The encoding of the \texttt{let} term in Twelf uses the
technique of \emph{higher-order abstract syntax} (HOAS)~\cite{Twelf-HOAS}.
HOAS is a technique for representing abstract syntax trees
with bound variables.
The abstract syntax from Section~\ref{sec:grammar} used to describe
the grammar of $\minilang$ is actually
\emph{first-order abstract syntax} (FOAS).
In FOAS, each AST has the form $o(t_1, t_2, \ldots, t_n)$,
where $o$ is an operator and $t_1, t_2, \ldots, t_n$ are each
AST themselves.
The operands of an operator correspond to subexpressions.
For example, in ``\texttt{+($e_1$; $e_2$)}'',
$e_1$ and $e_2$ are operands/subexpressions of the
\texttt{+} operator.
The Twelf encoding of ``\texttt{+($e_1$; $e_2$)}''
is ``\texttt{add : exp -> exp -> exp.}'',
where the \texttt{add} corresponds to \texttt{+},
the leftmost \texttt{exp} corresponds to $e_1$,
and the middle \texttt{exp} corresponds to $e_2$.

In HOAS, ASTs declare the variables they bind.
For example, consider the AST
$o(t_1, t_2, \ldots, t_n)$.
In HOAS, each $t_i$ has the form
$x_1, x_2, \ldots x_k.t$, where $t$ is an AST,
each $x_j$ is a variable bound in $t$, and $k \geq 0$;
if $k = 0$, then $t_i$ does not introduce new variables in $t$.
One advantage of knowing each variable introduced by an AST
is that we can rename variables without changing the meaning
the AST.

Now consider the \texttt{let} expression in $\minilang$ and
its representation in FOAS:
\[\code{let($x$; $e_1$; $e_2$)}\]
Recall that this \code{let} expression is evaluated by
binding the variable $x$ to the value of $e_1$ in $e_2$.
Hence, the \code{let} expression can be represented by the
following HOAS:
\[\code{let($e_1$; $x.e_2$)}\]
The ``$x.e_2$'' captures that $x$ is bound in $e_2$.
HOAS lets us know where variables are being bound.
This lets us easily determine that the following two expressions
are equivalent, since they only differ in variable names:
\[ \code{let(3; $x$.+($x$; 4))} \equiv \code{let(3; $y$.+($y$; 4))}\]

Our Twelf encoding of the \code{let} expression encodes its
HOAS version, which can be seen with the type of its
second argument: \code{(exp -> exp)}.
Usually, each operator such as \code{add} and \code{cat}
only took in \code{exp} arguments, where the arguments
represented subexpressions.
However, the second argument to \code{let} is not just an
ordinary subexpression but a subexpression that introduces
a new variable.
The second argument to \code{let} represents a higher-order
AST of the form ``$x.e_2$''.

An AST that introduces new variables is represented in
Twelf by a function.
Functions in programming languages are really just terms
with \emph{holes}.
The holes are represented by free variables,
and these holes are filled in when these
(terms w/ holes)/functions are applied to other terms.
Hence, ``$x.e_2$'' can be thought of as the function
or \emph{lambda-abstraction} ``$\lambda x.e_2$'',
where $e_2$ is the body of the function and $x$ is a variable
that is bound to a term in $e_2$.
A term $t$ of type \code{(exp -> exp)} will be a term of type
\code{exp} that contains occurrences of a free variable that
will be replaced with another term when $t$ is applied to it.
Twelf's syntax for the function ``$\lambda x.e$'' is
``\code{[$x$] $e$}''.
Figure~\ref{fig:hoas} gives a concrete example of a
\texttt{let} expression in
concrete syntax and its corresponding representation
in Twelf's HOAS.

\begin{figure}[H]
\begin{center}
\begin{tabular}{l|l}
\textit{Concrete Syntax} & \textit{Twelf HOAS} \\ \hline
\code{let x = $\underbrace{\code{1 + 2}}_{e_1}$ in
                $\underbrace{\code{x + 3}}_{e_2}$} &
  \code{let $\underbrace{\code{(add 1 2)}}_{e_1}$
    $\underbrace{\code{([x] add x 3)}}_{x.e_2}$}
\end{tabular}
\end{center}
\caption{Example Higher-Order Abstract Syntax in Twelf}
\label{fig:hoas} 
\end{figure}


